\chapter{Introduction}
\label{Introduction}

This section clarifies the circumstances and the motivation for the following research,
which problems are addressed, how we address them and what was achieved.

\section{Motivation}

The \ac{HPC} community and the \ac{CC} community have traditionally operated in separate worlds,
with the \ac{HPC} community focusing on raw performance and the \ac{CC} community focusing on convenience, flexibility, and reliability.
But with the rise of \ac{CC} systems to run a significant portion of the computation workloads of the Internet,
and a mature ecosystem that rivals the compute power of traditional \ac{HPC} systems,
both communities have realized the potential of combining the best of both worlds.

% Especially in the field of \ac{AI} and \ac{ML}, where increasingly large systems are needed for training large models,
% innovative solutions are being created to run these workloads directly on \ac{CC} systems.
% The recent acquisition of \texttt{Pachyderm} by \texttt{HPE} is an example of 
% one of the leading \ac{HPC} vendors sees the potential of combining the two worlds\footcite{nahrenkhizeranHewlettPackardEnterprise}.
% [Note from Harumi: Pachyderm is not an HPC system.]
% Correction:
Especially in the field of \ac{AI} and \ac{ML}, where increasingly large systems are needed for training large models,
innovative solutions are being created to run these workloads directly on \ac{CC} systems.
The recent acquisition of \texttt{Pachyderm} by \texttt{HPE} is an example of how one of the leading \ac{HPC} vendors
sees the potential benefits of \ac{CC} systems which can be used to schedule and run Big Data workloads in a
scalable and scheduled manner\footcite{nahrenkhizeranHewlettPackardEnterprise}.

% Within the \texttt{Hewlett Packard Labs - Systems Architecture Lab} a research project was started 
% to investigate the potential of just the workloads themselves but features of the \ac{CC} ecosystem to the \ac{HPC} world as well.
% [Note from Harumi: I think it is the other way round -- bringing CC turnkey convenience to HPC] 
% Correction: (Honeslty I think the original sentence was quite confusing in general)
Within the \texttt{Hewlett Packard Labs - Systems Architecture Lab} a research project was started in order
to investigate the potential of bringing the flexibility and convenience of the \ac{CC} ecosystem to the \ac{HPC} world.
While the \texttt{Pachykouda} project has successfully shown that the integration of the workflow management system \texttt{Pachyderm}
and the highly parallel compute framework \texttt{Arkouda} is possible, further research was needed.

The existing solution is based on a \gls{GitOps} approach, defining workflows, building containers and deploying them to the Kubernetes cluster,
based on the changes in a Git repository. While this approach is highly reliable and ensures reproducibility of the results,
it is limited by its static nature and a high latency between changes made and the deployment of the new workflow.
This latency prevents the system from being used in an interactive manner, as the user has to wait for the deployment to finish before the results can be seen.

The goal of this research project is not to replace the existing solution, but to extend the system with a new interactive way of
submitting jobs to the \ac{CC} system, allowing the user to interact with the system in real-time and get immediate feedback on the results,
without incurring the initial latency of the deployment, the build and deployment steps of the \gls{GitOps} approach.

\section{Problem Statement}
\label{ProblemStatement}

While low latency high performance computing has been a focus of the \ac{HPC} as well as the \ac{CC} community,
resulting in systems such as \texttt{JupyterHub} or \texttt{Dask}, these systems are inherently ephemeral and do not provide a way to reproduce results on a \ac{HPC} scale.
Therefore, we need to solve the following problems:

\begin{itemize}
    \item \textbf{P1:} \textit{The \gls{GitOps} first approach with introduces a latency between the changes made and the deployment of the new workflow.
              This latency prevents a real-time explorative workflow and limits the usability for the interactive \iffalse [Note from Harumi: do you mean interactive?] [Response from Jon: Yes :D] \fi exploration of the data.
          }
    \item \textbf{P2:} \textit{Existing solutions for interactive workloads are focusing on the convenience and flexibility of \ac{CC} systems,
              but do not provide a build in way to be combined with an automatic reproducible orchestration system for \ac{HPC} scale workloads.
          }
\end{itemize}

\section{Research Questions}
Based on these needs and problems the following questions need to be answered:

\begin{itemize}
    \item \textbf{RQ1:} \textit{
              How can the system be extended to allow for a real-time explorative workflow, allowing for an interactive exploration of \ac{HPC} scale data?
          }
    \item \textbf{RQ2:} \textit{
              How can such a system be designed to combine the performance of \ac{HPC} with the flexibility and convenience of \ac{CC} systems?
          }

\end{itemize}


\section{Contributions}

By way of this research the following results were achieved:

\begin{itemize}
    \item \textbf{C1:} \textit{ A deductive analysis of the three intersecting fields of \ac{HPC}, \ac{CC} and \ac{EC} }
    \item \textbf{C2:} \textit{ An abstraction of these fields into a conceptual solution for the problem at hand}
    \item \textbf{C3:} \textit{ A practical implementation of the conceptual solution showing the feasibility of the approach, where the deductive analysis and the conceptual solution are used as a guideline and are validated by the practical implementation.}
\end{itemize}

\pagebreak
